\section{问题重述}

\subsection{问题背景}

大语言模型（LLM）的能力由参数规模 $N$、训练数据量 $D$ 与数据质量 $Q$ 共同决定。
在算力资源受限的现实条件下，如何把有限算力在``扩大模型''、``增加数据''与
``提升数据质量''三条途径之间合理配置，使模型能力（以验证损失 $L$ 度量）最优，
是一个兼具理论与工程价值的问题。本题给定三个数据族：A 族为数据质量信号与训练
配比数据，B 族为缩放实验与损失轨迹数据，C 族为模型演进、算力与排行榜数据，
要求在真实数据上完成建模、求解与预测。

\subsection{问题内容}

\textbf{问题一：数据质量度量与领域配比建模。}（1）基于 A1--A3 的质量信号字段
（22 项指标，含标量与列表两类），建立可解释的数据质量评分方法，给出样本级与
域级质量分；（2）定义并消解质量信号之间的冲突，给出高冲突样本的处理规则；
（3）基于 A4--A15 的配比--损失数据，建立 17 个领域配比与各损失域之间的数量
关系模型，并讨论配比效应随模型规模的转移。

\textbf{问题二：含质量维度的广义标度律。}在经典标度律基础上，把数据质量 $Q$
纳入能力方程，建立广义标度律 $L(N,D,Q)$；利用 B 族数据估计参数，验证跨家族
泛化能力，并分析各维度的弹性与相互替代关系。

\textbf{问题三：算力约束下的联合优化与结构性转移。}在总算力预算
$C\in\{10^{19},10^{22},10^{24}\}$ FLOPs 下，给定训练成本
$C_{\text{train}}=6ND$、质量提升成本 $C_Q=D[g(Q)-g(Q_0)]_+$、长上下文注意力
成本 $C_{\text{attn}}=\eta N D L_{\text{ctx}}$（$\eta=2\times10^{-4}$，
$L_{\text{ctx}}$ 外生且取值须依据附件 C7），联合优化 $N$、$D$、$Q$（配比 $p$
可作为联立变量或由前两问确定）；分析预算跨越量级时最优策略是否发生``结构性
转移''，给出明确的数学定义与识别方法，并做灵敏度分析。

\textbf{问题四：技术演进分解与前沿预测。}基于附件 C，把开源模型能力提升分解为
规模扩张与非规模技术进步两部分并量化各自贡献占比；在算力增长放缓情景下，预测
未来 12/24 个月开源模型能力前沿，并给出不确定区间；同时完成至少一项基于 C8 的
逐任务聚合分析。

\subsection{已知条件与数据}

\begin{itemize}
\item A 族：A1（Slimpajama 质量信号样本 51230 条）、A2（arxiv 扩展 17523 条）、
A3（github 扩展 203752 条）、A4--A15（17 域配比与 13 个损失域的
1M/60M/1B/10b/70b 五档规模数据）、A16（领域映射表）、A18（可选用验算数据）；
\item B 族：B1（Pythia 训练轨迹 1176 条）、B2（Cerebras 训练日志 1029 条）、
B3（8 条完整轨迹）、B4/B5（缩放基准与文献数据）、B6--B8（$N$--$D$--$Q$
半合成实验，注意 B8 为外推补充）、B9/B10（大规模模型补充数据）；
\item C 族：C1/C2（开源模型排行榜 4576 条）、C3（扩展时序 4599 条）、
C4（Epoch 模型库 3523 条）、C5/C6（Loss--Benchmark 桥接 75 条）、
C7（模型架构元数据 45 条）、C8（逐任务明细 1863 个目录）。
\end{itemize}

\subsection{输出要求}

四问依次给出：质量评分方法与域级质量、广义标度律及其参数与验证、三档预算下
最优配置与结构性转移的数学刻画、前沿分解与 12/24 个月预测（含不确定性），
并附全部可复算代码与数据口径说明。
